\chapter{Resumen Ejecutivo}

El proyecto \textbf{Gestión de Comandas} consiste en el desarrollo de una aplicación
móvil integral diseñada para centralizar y gestionar operativamente los pedidos,
ventas, gastos y reportes financieros de restaurantes pequeños y medianos. La
plataforma aborda la problemática del registro manual en papel mediante una solución
tecnológica que permite a meseros y administradores operar en tiempo real desde
dispositivos Android, complementada con una carta digital pública y un sistema de
interacción directa del cliente desde su mesa.

\section{Problemática Resuelta}

La aplicación resuelve los principales desafíos identificados en la operación diaria
de restaurantes:

\begin{itemize}[leftmargin=1.5em]
  \item \textbf{Errores de cálculo} en el cierre diario de caja por registros manuales.
  \item \textbf{Dificultad de conciliación} entre ventas registradas e ingresos reales.
  \item \textbf{Ausencia de visibilidad financiera} en tiempo real para el administrador.
  \item \textbf{Procesos lentos} de verificación al final de cada turno.
  \item \textbf{Falta de trazabilidad} del estado de cada mesa durante el servicio.
  \item \textbf{Ausencia de un canal directo} entre cliente y mesero sin interrupciones físicas.
\end{itemize}

\section{Solución Implementada}

Se desarrolló una aplicación móvil nativa para Android con las siguientes
características principales:

\textbf{Arquitectura Técnica:}
\begin{itemize}[leftmargin=1.5em]
  \item \textit{Frontend móvil:} React Native con Expo SDK 52+
  \item \textit{Web (carta digital + sesión de cliente):} Expo Router (rutas públicas web)
  \item \textit{BaaS:} Supabase (PostgreSQL + Auth + Storage + Realtime)
  \item \textit{Lenguaje:} TypeScript
  \item \textit{Deploy:} Expo EAS Build + EAS Update (OTA)
\end{itemize}

\textbf{Funcionalidades Clave:}
\begin{itemize}[leftmargin=1.5em]
  \item Registro de pedidos con selección de platillos, mesa y método de pago.
  \item Vista en tiempo real del estado de mesas para el mesero.
  \item Notificaciones push cuando un pedido pasa a estado \textit{listo}.
  \item Carta digital pública accesible por QR sin necesidad de app ni registro.
  \item Sesión anónima del cliente por mesa: ver menú y llamar al mesero.
  \item Gestión completa del menú (CRUD + imágenes vía Supabase Storage).
  \item Registro de gastos operativos y generación de reportes financieros.
  \item Acceso diferenciado por rol: \texttt{waiter}, \texttt{admin}, \texttt{customer}.
\end{itemize}

\section{Metodología y Gestión}

El proyecto se desarrolló bajo la metodología ágil \textbf{Scrum}, organizado en
\textbf{3 sprints de 4 semanas} cada uno, con un equipo multidisciplinario de
5 integrantes que cumplieron roles específicos de desarrollo y gestión. Se utilizaron
herramientas modernas de colaboración como GitHub, Notion, Discord y Excalidraw
para garantizar una comunicación efectiva y un control de versiones robusto.